\documentclass[10pt,a4paper]{amsart}
\usepackage{amsmath,amssymb,amsthm,graphicx}
\usepackage{verbatim}
\newenvironment{reference}{\comment}{\endcomment}
%\newenvironment{reference}{}{}
\newenvironment{slogan}{\comment}{\endcomment}
\newenvironment{history}{\comment}{\endcomment}

% For commutative diagrams we use Xy-pic
\usepackage[all]{xy}

% We use 2cell for 2-commutative diagrams.
\xyoption{2cell}
% Theorem environments.
%
\theoremstyle{plain}
\newtheorem{theorem}[subsection]{Theorem}
\newtheorem{proposition}[subsection]{Proposition}
\newtheorem{lemma}[subsection]{Lemma}

\theoremstyle{definition}
\newtheorem{definition}[subsection]{Definition}
\newtheorem{example}[subsection]{Example}
\newtheorem{exercise}[subsection]{Exercise}
\newtheorem{situation}[subsection]{Situation}

\theoremstyle{remark}
\newtheorem{remark}[subsection]{Remark}
\newtheorem{remarks}[subsection]{Remarks}

\numberwithin{equation}{subsection}

% Macros
%
\def\lim{\mathop{\mathrm{lim}}\nolimits}
\def\colim{\mathop{\mathrm{colim}}\nolimits}
\def\Spec{\mathop{\mathrm{Spec}}}
\def\Hom{\mathop{\mathrm{Hom}}\nolimits}
\def\Ext{\mathop{\mathrm{Ext}}\nolimits}
\def\SheafHom{\mathop{\mathcal{H}\!\mathit{om}}\nolimits}
\def\SheafExt{\mathop{\mathcal{E}\!\mathit{xt}}\nolimits}
\def\Sch{\mathit{Sch}}
\def\Mor{\mathop{\mathrm{Mor}}\nolimits}
\def\Ob{\mathop{\mathrm{Ob}}\nolimits}
\def\Sh{\mathop{\mathit{Sh}}\nolimits}
\def\NL{\mathop{N\!L}\nolimits}
\def\CH{\mathop{\mathrm{CH}}\nolimits}
\def\proetale{{pro\text{-}\acute{e}tale}}
\def\etale{{\acute{e}tale}}
\def\QCoh{\mathit{QCoh}}
\def\Ker{\mathop{\mathrm{Ker}}}
\def\Im{\mathop{\mathrm{Im}}}
\def\Coker{\mathop{\mathrm{Coker}}}
\def\Coim{\mathop{\mathrm{Coim}}}

% Boxtimes
%
\DeclareMathSymbol{\boxtimes}{\mathbin}{AMSa}{"02}

%
% Macros for moduli stacks/spaces
%
\def\QCohstack{\mathcal{QC}\!\mathit{oh}}
\def\Cohstack{\mathcal{C}\!\mathit{oh}}
\def\Spacesstack{\mathcal{S}\!\mathit{paces}}
\def\Quotfunctor{\mathrm{Quot}}
\def\Hilbfunctor{\mathrm{Hilb}}
\def\Curvesstack{\mathcal{C}\!\mathit{urves}}
\def\Polarizedstack{\mathcal{P}\!\mathit{olarized}}
\def\Complexesstack{\mathcal{C}\!\mathit{omplexes}}
% \Pic is the operator that assigns to X its picard group, usage \Pic(X)
% \Picardstack_{X/B} denotes the Picard stack of X over B
% \Picardfunctor_{X/B} denotes the Picard functor of X over B
\def\Pic{\mathop{\mathrm{Pic}}\nolimits}
\def\Picardstack{\mathcal{P}\!\mathit{ic}}
\def\Picardfunctor{\mathrm{Pic}}
\def\Deformationcategory{\mathcal{D}\!\mathit{ef}}
\usepackage{cleveref}
\crefname{originalequation}{original equation}{original equations}
\Crefname{originalequation}{Original equation}{Original equations}
\makeatletter
\newlabel{equation-C-prime}{{0AL4}{source}{}{sourcepointer-1}{}}
\newlabel{equation-C-prime@cref}{{[originalequation][0][]0AL4}{[1][0][]source}{}{}{}}
\makeatother
\crefname{originallemma}{original lemma}{original lemmas}
\Crefname{originallemma}{Original lemma}{Original lemmas}
\makeatletter
\newlabel{local-cohomology-lemma-bound-two-term-complex}{{0GAD}{source}{}{sourcepointer-2}{}}
\newlabel{local-cohomology-lemma-bound-two-term-complex@cref}{{[originallemma][0][]0GAD}{[1][0][]source}{}{}{}}
\makeatother
\crefname{originalsection}{original section}{original sections}
\Crefname{originalsection}{Original section}{Original sections}
\makeatletter
\newlabel{local-cohomology-section-uniform}{{0GA6}{source}{}{sourcepointer-3}{}}
\newlabel{local-cohomology-section-uniform@cref}{{[originalsection][0][]0GA6}{[1][0][]source}{}{}{}}
\makeatother
\crefname{originallemma}{original lemma}{original lemmas}
\Crefname{originallemma}{Original lemma}{Original lemmas}
\makeatletter
\newlabel{more-algebra-lemma-two-term-base-change}{{0G9G}{source}{}{sourcepointer-4}{}}
\newlabel{more-algebra-lemma-two-term-base-change@cref}{{[originallemma][0][]0G9G}{[1][0][]source}{}{}{}}
\makeatother
\crefname{originallemma}{original lemma}{original lemmas}
\Crefname{originallemma}{Original lemma}{Original lemmas}
\makeatletter
\newlabel{more-algebra-lemma-upgrade-adjoint-tensor-RHom}{{0E1W}{source}{}{sourcepointer-5}{}}
\newlabel{more-algebra-lemma-upgrade-adjoint-tensor-RHom@cref}{{[originallemma][0][]0E1W}{[1][0][]source}{}{}{}}
\makeatother
\crefname{originalreference}{original reference}{original references}
\Crefname{originalreference}{Original reference}{Original references}
\makeatletter
\newlabel{remark-improve-homomorphism}{{0GAP}{source}{}{sourcepointer-6}{}}
\newlabel{remark-improve-homomorphism@cref}{{[originalreference][0][]0GAP}{[1][0][]source}{}{}{}}
\makeatother
\title{Lifting an algebra map B to C modulo a deeper power of I}
\author{The Stacks Project Authors}
\date{stacks-project:restricted.tex:local-snapshot-20261012; GNU FDL 1.2 or later}
\begin{document}
\maketitle
\noindent\textit{Editorial note.} All delivered mathematical text is the author's own, copied from the pinned Stacks Project chapter file restricted.tex; no statement or proof step was rewritten, reordered or removed.\par\medskip
\noindent\textit{Editorial note.} The counted proof is the author's complete argument: the proof computes the Ext identifications, applies the uniform bound to push the obstruction past I\textasciicircum{}\{n'\}, observes 2n' > n from the hypothesis on n, iterates the lifting, and takes the limit in the complete ring C\par\medskip
\noindent\textit{Editorial note.} The label lemma-get-morphism-general-better names the counted statement, and the proof environment immediately following it in the source is delivered verbatim as its proof.\par\medskip
\noindent\textit{Editorial note.} The counted statement is carried by the author's own lemma environment in the chapter source.\par\medskip
\noindent\textit{Allowed standard prerequisites (editorial).}\begin{itemize}
\item the obstruction class o(psi\_n) in Ext\textasciicircum{}1\_B(NL\textasciicircum{}wedge\_\{B/A\}, I\textasciicircum{}nC/I\textasciicircum{}\{2n\}C) attached to a map psi\_n : B -> C/I\textasciicircum{}nC, constructed in the same chapter
\item the adjunction identification of Ext after tensoring with C/I\textasciicircum{}nC and the two-term base change identification, both from More on Algebra
\item the uniform bound lemma of Local Cohomology: integers d and q\_0 of the associated graded Gr\_I(C) control the annihilation of Ext classes of two-term complexes
\item I-adic completeness of C and taking limits of compatible algebra maps
\end{itemize}\medskip
\ifdefined\hypertarget\hypertarget{sourcepointer-1}{}\fi
\noindent\textit{Source reference.} Original equation 0AL4: sources/stacks-restricted/pinned-original-source.tex, source lines 326--326. Prerequisite pointer invoked at proof line(s) 1172; it names a standard fact of the same chapter that the counted argument uses, and every step of the counted proof built on it is retained verbatim in the delivered body.\par\medskip
\ifdefined\hypertarget\hypertarget{sourcepointer-2}{}\fi
\noindent\textit{Source reference.} Original lemma 0GAD: sources/stacks-chow/tags.txt, source lines 20333--20333. Prerequisite pointer invoked at proof line(s) 1207; it names a standard fact of a sibling chapter that the counted argument uses, and every step of the counted proof built on it is retained verbatim in the delivered body.\par\medskip
\ifdefined\hypertarget\hypertarget{sourcepointer-3}{}\fi
\noindent\textit{Source reference.} Original section 0GA6: sources/stacks-chow/tags.txt, source lines 20326--20326. Prerequisite pointer invoked at proof line(s) 1177; it names a standard fact of a sibling chapter that the counted argument uses, and every step of the counted proof built on it is retained verbatim in the delivered body.\par\medskip
\ifdefined\hypertarget\hypertarget{sourcepointer-4}{}\fi
\noindent\textit{Source reference.} Original lemma 0G9G: sources/stacks-chow/tags.txt, source lines 20301--20301. Prerequisite pointer invoked at proof line(s) 1202; it names a standard fact of a sibling chapter that the counted argument uses, and every step of the counted proof built on it is retained verbatim in the delivered body.\par\medskip
\ifdefined\hypertarget\hypertarget{sourcepointer-5}{}\fi
\noindent\textit{Source reference.} Original lemma 0E1W: sources/stacks-chow/tags.txt, source lines 17586--17586. Prerequisite pointer invoked at proof line(s) 1200; it names a standard fact of a sibling chapter that the counted argument uses, and every step of the counted proof built on it is retained verbatim in the delivered body.\par\medskip
\ifdefined\hypertarget\hypertarget{sourcepointer-6}{}\fi
\noindent\textit{Source reference.} Original reference 0GAP: sources/stacks-restricted/pinned-original-source.tex, source lines 1101--1101. Prerequisite pointer invoked at proof line(s) 1190; it names a standard fact of the same chapter that the counted argument uses, and every step of the counted proof built on it is retained verbatim in the delivered body.\par\medskip
\begin{lemma}
\label{lemma-get-morphism-general-better}
Assume given the following data
\begin{enumerate}
\item an integer $c \geq 0$,
\item an ideal $I$ of a Noetherian ring $A$,
\item $B$ in (\ref{equation-C-prime}) for $(A, I)$ such that
$I^c$ annihilates $\Ext^1_B(\NL_{B/A}^\wedge, N)$
for any $B$-module $N$,
\item a Noetherian $I$-adically complete $A$-algebra $C$; denote
$d = d(\text{Gr}_I(C))$ and $q_0 = q(\text{Gr}_I(C))$ the integers found in
Local Cohomology, Section \ref{local-cohomology-section-uniform},
\item an integer $n \geq \max(q_0 + (d + 1)c, 2(d + 1)c + 1)$, and
\item an $A$-algebra homomorphism $\psi_n : B \to C/I^nC$.
\end{enumerate}
Then there exists a map $\varphi : B \to C$ of $A$-algebras such
that $\psi_n \bmod I^{n - (d + 1)c} = \varphi \bmod I^{n - (d + 1)c}$.
\end{lemma}

\begin{proof}
Consider the obstruction class
$$
o(\psi_n) \in \Ext^1_B(\NL_{B/A}^\wedge, I^nC/I^{2n}C)
$$
of Remark \ref{remark-improve-homomorphism}. For any $C/I^nC$-module
$N$ we have
\begin{align*}
\Ext^1_B(\NL_{B/A}^\wedge, N)
& =
\Ext^1_{C/I^nC}(\NL_{B/A}^\wedge \otimes_B^\mathbf{L} C/I^nC, N) \\
& =
\Ext^1_{C/I^nC}(\NL_{B/A}^\wedge \otimes_B C/I^nC, N)
\end{align*}
The first equality by
More on Algebra, Lemma \ref{more-algebra-lemma-upgrade-adjoint-tensor-RHom}
and the second one by
More on Algebra, Lemma \ref{more-algebra-lemma-two-term-base-change}.
In particular, we see that
$\Ext^1_{C/I^nC}(\NL_{B/A}^\wedge \otimes_B C/I^nC, N)$ is annihilated by
$I^cC$ for all $C/I^nC$-modules $N$.
It follows that we may apply
Local Cohomology, Lemma \ref{local-cohomology-lemma-bound-two-term-complex}
to see that $o(\psi_n)$ maps to zero in
$$
\Ext^1_{C/I^nC}(\NL_{B/A}^\wedge \otimes_B C/I^nC, I^{n'}C/I^{2n'}C) =
\Ext^1_B(\NL_{B/A}^\wedge, I^{n'}C/I^{2n'}C) =
$$
where $n' = n - (d + 1)c$. By the discussion in
Remark \ref{remark-improve-homomorphism} we obtain a map
$$
\psi'_{2n'} : B \to C/I^{2n'}C
$$
which agrees with $\psi_n$ modulo $I^{n'}$.
Observe that $2n' > n$ because $n \geq 2(d + 1)c + 1$.

\medskip\noindent
We may repeat this procedure. Starting with $n_0 = n$ and
$\psi^0 = \psi_n$ we end up getting a strictly increasing
sequence of integers
$$
n_0 < n_1 < n_2 < \ldots
$$
and $A$-algebra homorphisms $\psi^i : B \to C/I^{n_i}C$
such that $\psi^{i + 1}$ and $\psi^i$ agree modulo $I^{n_i - tc}$.
Since $C$ is $I$-adically complete we can take $\varphi$
to be the limit of the maps
$\psi^i \bmod I^{n_i - (d + 1)c} : B \to C/I^{n_i - (d + 1)c}C$
and the lemma follows.
\end{proof}

\end{document}
